\section{Financial Services on Astra}
\sectionkicker{FINANCIAL VERTICAL}
\begin{wideflow}
{\Large\bfseries 金融業務向けChatGPT、Astraの推論を中核に}\par
\smallskip
{\color{SurveyMuted}9月10日発表の金融業務向けパッケージを、内蔵データと評価・ガバナンスの留保と切り分けて示す。}\par
\end{wideflow}

9月10日、OpenAIはChatGPT for Financial Servicesを発表した\cite{finserv}。ChatGPT Workの操作感を土台に、GPT-6 Astraの推論能力をリサーチや財務モデリング、顧客向け資料作成に結びつける内容で、Morgan StanleyとEvercoreとの設計パートナーシップで磨いたとされる。当面の対象は投資銀行業務と株式リサーチである。内蔵データとしてDaloopa、PitchBook、LSEG News、Crunchbaseが挙がり、提供元の紹介ではFiscal.aiも添えられる。S\&P Capital IQやLSEG、MSCI、Dow Jones Factiva、Moody'sとのシングルサインオン連携、50を超えるコネクターが記される。管理側ではExcel・Word・PowerPointテンプレートの配布、SAMLやSCIM、ロールベースのアクセス制御、暗号化や保存期間の設定などエンタープライズ向けのガバナンスが示される。価格や利用枠、提供地域は資料に見当たらず、本号では断定しない。データ精度や応答速度の改善はベンダー側の説明として受け止める。

評価として掲げられるのはOfficeQA Proの69.9\%で、GPT-5.6 Solの60.2\%と並べられる\cite{finserv}。米国債ブレティンなど表・グラフ・脚注をまたぐ情報探索を問う設計と説明されるが、OfficeQAやBoxBench、スライド作成タスクの採点の詳細は今回たどっておらず、サンプリングや判定の妥当性は確認していない。いずれもベンダー側の数値であり、独立した再現を待つ。Astra自体の位置づけも、検索・推論・成果物生成に強いというベンダー側の表現として扱い、金融向けの一般的な優劣として広げない。この節の内容をGPT-Live-1やAgents APIの仕様として読み替えない。3つは独立したリリースである。

\begin{claimboundary}[資料と評価の境界]
価格、利用枠、提供地域は資料にない。評価の詳細はたどっていない。数値はベンダー側の報告であり、独立した再現を待つ。
\end{claimboundary}

\begin{communitynote}[X / Community observation --- context only]
公開投稿では、9月10日の公式発信に独立した言及が重なる様子が観察される\cite{w37x-astra-official,w37x-astra-ind}。うかがえるのは利用枠への関心や使い比べへの注目であり、仕様やベンチマーク、料金、利用条件の根拠にはならない。締め切り後の投稿は文脈資料であり、通常の集計には含めない。
\end{communitynote}
